\begin{textsample}{NEG Dim 1 – human – Score -6.00 – mg/mg\_ef\_linguagens\_ed\_fis\_5\_p1.txt}  \label{ex:f1_neg_007}
Esse princípio se funda na premissa de que o conhecimento sobre o corpo e vivido no corpo é que nos \textbf{possibilita} compreender a nossa existência no mundo, pois é por meio dele que \textbf{construímos} \textbf{significados}, ocupamos espaços, comunicamos, interagimos e nos constituímos como identidades individuais e \textbf{coletivas}. É, portanto, com base nesse pressuposto que concebemos as práticas \textbf{corporais} como linguagem.
Várias são as concepções de linguagem.
Entretanto, como queremos que nossos estudantes sejam capazes de ler, interpretar e \textbf{produzir} diversos tipos de textos - gestuais, orais, escritos, virtuais e outros – com senso crítico, argumentativo, de modo a compreender os limites e as possibilidades de sua vivência social, entendemos que a concepção de linguagem como enunciação constitutiva é um caminho importante.

% matched lemmas: coletivo, construir, corporal, possibilitar, produzir, significado
\end{textsample}
